% !TEX root = NotasCursoOptimizacion.tex

%===========================================
\section{Programaci\'on Lineal}
%===========================================

\section{Graficar}


\begin{Ejer}
\begin{multicols}{3}
\begin{enumerate}
\item $\left\{
\begin{array}{c}
2x+5y\leq25\\
6x+5y\leq 45\\
x,y\geq0\end{array}
\right.$

\item $\left\{
\begin{array}{c}
2x+3y\geq120\\
2x+y\geq80\\
x,y\geq0\end{array}
\right.$

\item $\left\{
\begin{array}{c}
x+2y\leq40\\
6x+5y\leq150\\
x,y\geq0\end{array}
\right.$

\item $\left\{
\begin{array}{c}
40x+60y\geq\\
30x+15y\geq180\\
x,y\geq0\end{array}
\right.$

\item $\left\{
\begin{array}{c}
0.8x+0.5y\leq60\\
0.3x+0.4y\leq50\\
x,y\geq0\end{array}
\right.$

\item $\left\{
\begin{array}{c}
50x+60y\geq2500\\
100x+60y\geq3000\\
x,y\geq0\end{array}
\right.$

\item $\left\{
\begin{array}{c}
2x+3y+2z\leq4400\\
4x+0y+3z\leq470\\
2x+5y+0z\leq430\end{array}
\right.$

\item $\left\{
\begin{array}{c}
x+2y+2z\leq2\\
2x+y+z\leq2\\
x+5y+z\leq3\\
x,y,z\geq0\end{array}
\right.$

\item $\left\{
\begin{array}{c}
2x+t\leq100\\
x+y\leq80\\
x\leq40\\
x,y\geq0\end{array}
\right.$

\item $\left\{
\begin{array}{c}
2x+3y\geq-3\\
2x-y-9\leq0\\
2x-5y-5\geq0\\
x,y\geq0\end{array}
\right.$

\item $\left\{
\begin{array}{c}
x+2y\leq12\\
2x+y\geq4\\
x-2y\leq6\\
x-y\geq0\\
x,y\geq0\end{array}
\right.$

\item $\left\{
\begin{array}{c}
x+3y\geq50\\
9x-8y\geq0\\
0\leq x\leq8\\
0\leq y\leq2\end{array}
\right.$


\item $\left\{
\begin{array}{c}
x+y\geq5\\
y\geq x+3\\
3x-y\geq1\\
y+2x\leq16\\
4y-x\leq22
\end{array}
\right.$


\item $\left\{
\begin{array}{c}
x+y\geq14\\
2x+3y\geq36\\
4x+y\geq16\\
x-3y\geq0\\
x,y\geq0
\end{array}
\right.$


\item \item $\left\{
\begin{array}{c}
4x+5y\geq20\\
7x+2y\geq14\\
x\leq y\\
x,y\geq0
\end{array}
\right.$

\item \item $\left\{
\begin{array}{c}
2x+y\leq10\\
x+3y\geq12\\
0\leq x\leq8\\
o\leq y\leq2
\end{array}\right.$
\end{enumerate}
\end{multicols}
\end{Ejer}

\section{Definiciones}

A linear programming problem (LP) is an optimization problem where:

\begin{enumerate}
\item We maximize or minimize a linear function of decision variables (objective function).
\item The decision variables satisfy a set of linear constraints.
\item Each variable may have sign restrictions ($x_j \ge 0$ or unrestricted).
\end{enumerate}

\subsection*{Feasible Region and Optimal Solution}

\begin{Def}  
The feasible region of an LP is the set of all points that satisfy all constraints and sign restrictions.
\end{Def}

\begin{Def}  
For a maximization problem, an optimal solution is a feasible point with the largest objective value.  
For a minimization problem, it is the feasible point with the smallest objective value.
\end{Def}

\subsection*{Canonical Linear Programming Problem}

\begin{Def}[Canonical Form]

\begin{enumerate}
\item $x_j \ge 0$, $j = 1,\dots,N$
\item $Ax = b$
\item Minimize $c^T x$
\end{enumerate}

Matrix form:

\[
\min c^T x
\]
\[
Ax = b
\]
\[
x \ge 0
\]
\end{Def}

\subsection*{Standard Maximum Problem}

Find $x \in \mathbb{R}^n$ to maximize

\[
c^T x = c_1x_1 + \cdots + c_n x_n
\]

subject to

\[
Ax \le b
\]
\[
x \ge 0
\]

where

\[
A =
\begin{pmatrix}
a_{11} & \cdots & a_{1n} \\
\vdots & \ddots & \vdots \\
a_{m1} & \cdots & a_{mn}
\end{pmatrix}
\]

\subsection*{Standard Minimum Problem}

Find $y \in \mathbb{R}^m$ to minimize

\[
y^T b
\]

subject to

\[
y^T A \ge c^T
\]
\[
y \ge 0
\]


\subsection*{Extreme Points and Optimality}

\begin{Obs}  
The optimal solution occurs at an extreme point of the feasible region.
\end{Obs}

\begin{Obs}  
A point is optimal if and only if the negative gradient of the objective function lies in the cone spanned by the gradients of the active constraints.
\end{Obs}

\subsection*{Resource Allocation Model}

Let:

\[
\max \sum_{j=1}^n c_j x_j
\]

subject to

\[
\sum_{j=1}^n a_{ij} x_j \le b_i, \quad i = 1,\dots,m
\]

\[
x_j \ge 0
\]

Matrix form:

\[
\max z = c^T x
\]
\[
Ax \le b
\]
\[
x \ge 0
\]


\subsection*{Basic Terminology}

\begin{Def}  
If $x$ satisfies $Ax=b$, $x\ge 0$, then $x$ is \emph{feasible}.
\end{Def}

\begin{Def}
An LP is feasible if it has a feasible solution.
\end{Def}


\begin{Def}
$x^*$ is optimal if

\[
c^T x^* = \min \{c^T x : Ax=b, x\ge 0\}.
\]
\end{Def}

\begin{Def}
An LP is unbounded if the objective function can decrease without bound over feasible solutions.
\end{Def}

\subsection*{Equivalent Forms}

\begin{enumerate}

\item Maximization as minimization:

\[
\max c^T x \quad \Longleftrightarrow \quad \min -c^T x.
\]

\item Equality as inequalities:

\[
a_i^T x = b_i
\quad \Longleftrightarrow \quad
\begin{cases}
a_i^T x \le b_i, \\
-a_i^T x \le -b_i.
\end{cases}
\tag{LP-2}
\]

\item Slack variables:

\[
a_i^T x \le b_i
\quad \Longleftrightarrow \quad
a_i^T x + s_i = b_i,
\quad s_i \ge 0.
\]

\item Non-positivity:

If $x_j \le 0$, replace $x_j = -y_j$, $y_j \ge 0$.

\item Unrestricted variable:

\[
x_j = x_j^+ - x_j^-,
\quad x_j^+, x_j^- \ge 0.
\]
\end{enumerate}

\subsection*{Geometry of Linear Programming}

Let

\[
P = \{x : Ax=b, x \ge 0\}.
\]

\begin{Def}
$x$ is a vertex of $P$ if there does not exist $y \neq 0$ such that
$x+y \in P$ and $x-y \in P$.
\end{Def}

\begin{Teo}  
If $\min \{c^T x : x\in P\}$ is finite, then there exists a vertex solution $x^*$ with
$c^T x^* \le c^T x$ for all $x \in P$.
\end{Teo}


\begin{Coro}
If $\min\{c^T x : Ax=b, x\ge 0\}$ is finite, then there exists an optimal solution which is a vertex.
\end{Coro}

\begin{Coro}
If $P=\{x:Ax=b, x\ge 0\}\neq \emptyset$, then $P$ has a vertex.
\end{Coro}


\subsection*{Standard Maximum and Minimum Problems}

\textbf{Standard Maximum}

\[
\max c^T x
\]

subject to

\[
Ax \le b,
\quad
x \ge 0.
\]

\textbf{Standard Minimum}

\[
\min y^T b
\]

subject to

\[
y^T A \ge c^T,
\quad
y \ge 0.
\]


\subsection{What is a Linear Programming (LP) Problem?}

\begin{Def}[Linear Function]
A function $f(x_1,x_2,\ldots,x_n)$ of the variables 
$x_1,x_2,\ldots,x_n$ is a \textbf{linear function} if

\[
f(x_1,x_2,\ldots,x_n)=c_1x_1+c_2x_2+\cdots+c_nx_n
\]

where $c_1,c_2,\ldots,c_n$ are constants.
\end{Def}

\begin{Ejem}
$f(x_1,x_2)=2x_1+x_2$ is linear, while 
$f(x_1,x_2)=x_1^2x_2$ is not linear.
\end{Ejem}

\begin{Def}[Linear Inequalities]
For any linear function $f(x_1,x_2,\ldots,x_n)$ and any real number $b$,
the inequalities

\[
f(x_1,x_2,\ldots,x_n)\le b
\quad \text{and} \quad
f(x_1,x_2,\ldots,x_n)\ge b
\]

are called \textbf{linear inequalities}.
\end{Def}

\begin{Ejem}
$2x_1+3x_2\le 3$ is a linear inequality, while 
$x_1x_2+\sqrt{x_2}\ge 3$ is not linear.
\end{Ejem}

\begin{Def}[Linear Programming Problem]
A \textbf{linear programming (LP) problem} is an optimization problem in which:

\begin{enumerate}
\item We attempt to maximize (profit) or minimize (cost) a linear function
(called the \textit{objective function}) of the decision variables.

\item The decision variables must satisfy a set of constraints, and each
constraint must be a linear equation or linear inequality.

\item A \textbf{sign restriction} is associated with each variable.
For any variable $x_i$, either $x_i \ge 0$ or $x_i$ is unrestricted in sign.
\end{enumerate}

The general structure is:

\[
\max \text{ (or min)} \quad z=f(x_1,x_2,\ldots,x_n)
\]

\[
\text{s.t.} \quad \text{Constraints}
\]

\[
\text{Sign restrictions}
\]
\end{Def}


\subsection*{LP Assumptions}

\textbf{1. Proportionality and Additivity}

The objective function and constraints must be linear.

\begin{itemize}
\item Contributions are proportional to variable values.
\item Contributions are independent of other variables.
\end{itemize}

\textbf{2. Divisibility Assumption}

Decision variables can take fractional values.

If variables must be integers, the problem becomes an
\textbf{integer programming problem}.

\textbf{3. Certainty Assumption}

All parameters (objective coefficients, RHS values, constraint coefficients)
are known with certainty.


\subsection*{Linear Program in Matrix Form}

Standard form:

\[
\min \sum_{j=1}^{n} c_j x_j
\]

subject to

\[
\sum_{j=1}^{n} a_{ij} x_j = b_i,
\quad i=1,\dots,m,
\]

\[
x_j \ge 0.
\]

Matrix representation:

\[
\min c^T x
\quad \text{subject to}
\quad
Ax = b,
\quad
x \ge 0.
\tag{LP-1}
\]

Where

\[
x =
\begin{pmatrix}
x_1 \\ \vdots \\ x_n
\end{pmatrix},
\quad
b =
\begin{pmatrix}
b_1 \\ \vdots \\ b_m
\end{pmatrix},
\quad
c =
\begin{pmatrix}
c_1 \\ \vdots \\ c_n
\end{pmatrix},
\]

\[
A =
\begin{pmatrix}
a_{11} & \dots & a_{1n} \\
\vdots & \ddots & \vdots \\
a_{m1} & \dots & a_{mn}
\end{pmatrix}.
\]

\newpage

\section{Ejemplos}

\begin{Ejem}[Compañía de Pinturas]

La compañía de pinturas \textbf{Manchita} produce tres tipos de pintura adicionando a una pintura base cuatro productos químicos ($Q_1$ a $Q_4$). 

Se desea determinar la cantidad de toneladas de cada tipo de pintura que debe producirse para maximizar la ganancia total.

\bigskip

\begin{center}
\begin{tabular}{c|ccc|c}
 & \multicolumn{3}{c|}{Kg de aditivo por tonelada} & Disponible \\
Aditivo & Interior & Exterior & Especial & (kg) \\

$Q_1$ & 1 & 2 & 2 & 2 \\
$Q_2$ & 2 & 1 & 1 & 1 \\
$Q_3$ & 1 & 5 & 1 & 3 \\
$Q_4$ & 0 & 0 & 1 & 0.8 \\
Ganancia & 15000 & 25000 & 19000 & \\

\end{tabular}
\end{center}
\end{Ejem}

\begin{Ejem}[Producción en Máquinas]

Una empresa produce tres productos procesados en tres máquinas.

\begin{center}
\begin{tabular}{c|ccc|c}

 & \multicolumn{3}{c|}{Tiempo por unidad (min)} & Capacidad \\
Máquina & Prod.1 & Prod.2 & Prod.3 & (min/día) \\

$M_1$ & 2 & 3 & 2 & 440 \\
$M_2$ & 4 & - & 3 & 470 \\
$M_3$ & 2 & 5 & - & 430 \\

\end{tabular}
\end{center}
\end{Ejem}


\begin{Ejem}[Mezcla de Café]

Una casa empacadora recibe diariamente 700 kg de café tipo C y 800 kg tipo K.

Mezcla A:
\[
2 \text{ partes C} + 1 \text{ parte K}
\]
Ganancia: 22 céntimos/kg.

Mezcla B:
\[
1 \text{ parte C} + 2 \text{ partes K}
\]
Ganancia: 26 céntimos/kg.

Determinar las cantidades que maximizan la ganancia.
\end{Ejem}

\begin{Ejem}[Trenes y Camiones]

Un artesano fabrica trenes y camiones.

\begin{center}
\begin{tabular}{c|ccc}

 & Tornillos & Bloques & Ruedas \\

Tren & 10 & 15 & 18 \\
Camión & 20 & 10 & 6 \\

Disponibles & 8000 & 6000 & 6300 \\

\end{tabular}
\end{center}

Beneficios:
\[
1.6 \text{ €/tren}, \quad 1.4 \text{ €/camión}
\]
\end{Ejem}


\begin{Ejem}[Problema 5: Calentadores y Aires]

Producción diaria con restricciones de horas:

\begin{itemize}
\item Calentador: 2h partes + 1h ensamblaje
\item Aire acondicionado: 1h partes + 2h ensamblaje
\end{itemize}

Disponibilidad:
\[
8 \text{h partes}, \quad 10 \text{h ensamblaje}
\]

Ganancias:
\[
30 \text{ por calentador}, \quad 50 \text{ por aire}
\]
\end{Ejem}

\begin{Ejem}[Vino y Vinagre]

Restricciones:

\[
2V \leq G + 4
\]
\[
3G + 4V \leq 18
\]

Beneficio:
\[
8V + 2G
\]
\end{Ejem}


\begin{Ejem}[Tapicería]

Tipo A:
\[
10m \text{ tela}, \quad 12h
\]

Tipo B:
\[
15m \text{ tela}, \quad 16h
\]

Disponibilidad:
\[
300m \text{ tela}, \quad 336h
\]

Ganancias:
\[
1500A + 2100B
\]

\end{Ejem}


\begin{Ejem}[Gadgets y Gewgaws]

Modelo:

\[
\max z = 5x_1 + 4x_2
\]

Sujeto a:

\[
x_1 + 2x_2 \leq 120
\]
\[
x_1 + x_2 \leq 70
\]
\[
2x_1 + x_2 \leq 100
\]
\[
x_1, x_2 \geq 0
\]

\end{Ejem}


\begin{Ejem}[Modelo Simple]

\[
\max 3x_1 + x_2
\]

Sujeto a:

\[
x_1 + x_2 \leq 7
\]
\[
x_1 \geq 0
\]
\[
x_2 \geq 0
\]

\end{Ejem}


\begin{Ejem}[Problema 10: Dieta de Cerdos]

\begin{center}
\begin{tabular}{c|ccc|c}

 & Corn & Silage & Alfalfa & Requerimiento \\

Carbs & 0.9 & 0.2 & 0.4 & 2 \\
Protein & 3 & 8 & 6 & 18 \\
Vitamins & 1 & 2 & 4 & 15 \\
Cost & 7 & 6 & 5 & \\

\end{tabular}
\end{center}

\end{Ejem}


\begin{Ejem}[Modelo Dieta G1 y G2]

\begin{center}
\begin{tabular}{c|ccc|c}

 & Starch & Proteins & Vitamins & Cost \\

G1 & 5 & 4 & 2 & 0.6 \\
G2 & 7 & 2 & 1 & 0.35 \\

\end{tabular}
\end{center}

Requerimientos diarios:
\[
8 \text{ starch}, \quad 15 \text{ proteins}, \quad 3 \text{ vitamins}
\]

\end{Ejem}


\begin{Ejem}[Problema 12: Giapetto]

Modelo:

\[
\max Z = \text{Ingresos} - \text{Costos}
\]

Restricciones:
\[
2F + C \leq 100
\]
\[
F + C \leq 80
\]
\[
C \leq 40
\]

\end{Ejem}


\begin{Ejem}[Problema 13: Industria de Dos Productos]

\[
\max \text{Profit} = 3x_1 + 4x_2
\]

Sujeto a:

\[
2x_1 + 3x_2 \leq 16
\]
\[
4x_1 + 2x_2 \leq 16
\]
\[
x_1, x_2 \geq 0
\]

\end{Ejem}


\begin{Ejem}[Department Time Constraints]

Time used in making $x_1$ and $x_2$ units of two products must not exceed the total daily time available in the two shops.

\[
2x_1 + 3x_2 \le 16 \qquad \text{(Machine shop constraint)}
\]
\[
4x_1 + 2x_2 \le 16 \qquad \text{(Assembly shop constraint)}
\]

Nonnegativity constraints:
\[
x_1 \ge 0, \quad x_2 \ge 0.
\]

The problem in mathematical form:

\[
\textbf{Maximize} \quad Z = 3x_1 + 4x_2
\]

subject to
\[
2x_1 + 3x_2 \le 16
\]
\[
4x_1 + 2x_2 \le 16
\]
\[
x_1, x_2 \ge 0.
\]

\end{Ejem}


\begin{Ejem}[Graphical Solution]

Consider

\[
\max Z = x + 2y
\]

subject to

\[
x \ge 0, \quad y \ge 0,
\]
\[
x \le 6, \quad y \le 7,
\]
\[
x + y \le 12.
\]

\end{Ejem}


\begin{Ejem}[The Diet Problem]

There are $m$ different types of food $F_1,\dots,F_m$ that supply $n$ nutrients
$N_1,\dots,N_n$.

Let:
\begin{itemize}
\item $b_i$ = price per unit of food $F_i$,
\item $a_{ij}$ = amount of nutrient $N_j$ in one unit of $F_i$,
\item $c_j$ = minimum daily requirement of nutrient $N_j$,
\item $y_i$ = units of food $F_i$ purchased per day.
\end{itemize}

Cost per day:

\[
b_1 y_1 + b_2 y_2 + \dots + b_m y_m
\tag{1}
\]

Nutrient requirement:

\[
a_{1j}y_1 + a_{2j}y_2 + \dots + a_{mj}y_m \ge c_j,
\quad j = 1,\dots,n.
\tag{2}
\]

Nonnegativity:

\[
y_i \ge 0, \quad i=1,\dots,m.
\tag{3}
\]

Thus, minimize (1) subject to (2)–(3).
\end{Ejem}


\begin{Ejem}[The Transportation Problem]

There are $I$ production plants and $J$ markets.

Let:
\begin{itemize}
\item $s_i$ = supply at plant $P_i$,
\item $r_j$ = demand at market $M_j$,
\item $b_{ij}$ = transportation cost from $P_i$ to $M_j$,
\item $y_{ij}$ = amount shipped from $P_i$ to $M_j$.
\end{itemize}

Total transportation cost:

\[
\sum_{i=1}^{I} \sum_{j=1}^{J} y_{ij} b_{ij}.
\tag{4}
\]

Supply constraints:

\[
\sum_{j=1}^{J} y_{ij} \le s_i,
\quad i=1,\dots,I.
\tag{5}
\]

Demand constraints:

\[
\sum_{i=1}^{I} y_{ij} \ge r_j,
\quad j=1,\dots,J.
\tag{6}
\]

Nonnegativity:

\[
y_{ij} \ge 0.
\tag{7}
\]
\end{Ejem}


\begin{Ejem}[Compañía de Pinturas Manchita]

La compañía de pinturas \textbf{Manchita} produce tres tipos de pintura adicionando a una pintura base cuatro productos o aditivos químicos ($Q_1$ a $Q_4$). Se tiene abundante pintura base disponible y su costo ya fue cubierto.

La compañía desea determinar la cantidad de toneladas de cada tipo de pintura que debe producir de manera que maximice la ganancia total. Las únicas restricciones se deben a la disponibilidad de los aditivos químicos requeridos.

\begin{center}
\begin{tabular}{|c|c|c|c|c|}
\hline
\textbf{Aditivo} & \multicolumn{3}{c|}{\textbf{Kg requeridos por tonelada}} & \textbf{Disponible (kg)} \\
\cline{2-4}
 & Pintura interior & Pintura exterior & Pintura especial & \\
\hline
$Q_1$ & 1 & 2 & 2 & 2 \\
$Q_2$ & 2 & 1 & 1 & 1 \\
$Q_3$ & 1 & 5 & 1 & 3 \\
$Q_4$ & 0 & 0 & 1 & 0.8 \\
\hline
\textbf{Ganancia por tonelada} & 15000 & 25000 & 19000 & \\
\hline
\end{tabular}
\end{center}

\end{Ejem}


\begin{Ejem}[Producción en Tres Máquinas]

Una empresa produce tres productos procesados en tres máquinas diferentes. El tiempo requerido para fabricar una unidad de cada producto y la capacidad diaria de las máquinas se muestran en la tabla:

\begin{center}
\begin{tabular}{|c|c|c|c|c|}
\hline
\textbf{Máquina} & Producto 1 & Producto 2 & Producto 3 & \textbf{Capacidad (min/día)} \\
\hline
$M_1$ & 2 & 3 & 2 & 440 \\
$M_2$ & 4 & - & 3 & 470 \\
$M_3$ & 2 & 5 & - & 430 \\
\hline
\end{tabular}
\end{center}
\end{Ejem}


\begin{Ejem}[Mezclas de Café]

Una casa empacadora recibe diariamente 700 kg de café tipo C y 800 kg de café tipo K. Produce dos mezclas:

\begin{itemize}
\item Tipo A: 2 partes de café tipo C y 1 parte de café tipo K. Ganancia de 22 céntimos por kilo.
\item Tipo B: 1 parte de café tipo C y 2 partes de café tipo K. Ganancia de 26 céntimos por kilo.
\end{itemize}

Determinar la cantidad de cada mezcla que debe prepararse para maximizar la ganancia.
\end{Ejem}


\begin{Ejem}[Trenes y Camiones]

Un artesano fabrica trenes y camiones de juguete utilizando tornillos, bloques y ruedas.

Disponibilidad semanal:
\[
8000 \text{ tornillos}, \quad 6000 \text{ bloques}, \quad 6300 \text{ ruedas}
\]

Beneficios:
\[
\text{Tren: } 1.6 \text{ euros/unidad} \qquad
\text{Camión: } 1.4 \text{ euros/unidad}
\]

\begin{center}
\begin{tabular}{|c|c|c|c|}
\hline
 & Tornillos & Bloques & Ruedas \\
\hline
Tren & 10 & 15 & 18 \\
Camión & 20 & 10 & 6 \\
\hline
\end{tabular}
\end{center}

\end{Ejem}


\begin{Ejem}[Heaters and Air Conditioners]

An appliance company manufactures heaters and air conditioners.

\begin{itemize}
\item Heater: 2 hours (parts division), 1 hour (assembly division)
\item Air conditioner: 1 hour (parts division), 2 hours (assembly division)
\end{itemize}

Disponibilidad:
\[
\text{Parts division: } 8 \text{ hours/day}
\]
\[
\text{Assembly division: } 10 \text{ hours/day}
\]

Profit:
\[
\$30 \text{ per heater}, \quad \$50 \text{ per air conditioner}
\]

\end{Ejem}


\begin{Ejem}[Industria Vinícola]

Una industria produce vino y vinagre.

Restricciones:
\[
2(\text{vino}) \leq (\text{vinagre}) + 4
\]
\[
3(\text{vinagre}) + 4(\text{vino}) \leq 18
\]

Beneficio:
\[
8€ \text{ por vino}, \quad 2€ \text{ por vinagre}
\]
\end{Ejem}


\begin{Ejem}[Tapicería]

Tipos de sillones: A y B.

Requerimientos:

\begin{itemize}
\item Tipo A: 10 m de tela, 12 horas
\item Tipo B: 15 m de tela, 16 horas
\end{itemize}

Disponibilidad semanal:
\[
300 \text{ m tela}, \quad 336 \text{ horas}
\]

Utilidad:
\[
\$1500 \text{ (A)}, \quad \$2100 \text{ (B)}
\]
\end{Ejem}


\begin{Ejem}[Gadgets and Gewgaws]

Maximizar:
\[
\max z = 5x_1 + 4x_2
\]

sujeto a:

\[
x_1 + 2x_2 \leq 120
\]
\[
x_1 + x_2 \leq 70
\]
\[
2x_1 + x_2 \leq 100
\]
\[
x_1, x_2 \geq 0
\]
\end{Ejem}


\begin{Ejem}[Ejemplo sencillo]

Maximizar:
\[
3x_1 + x_2
\]

Sujeto a:
\[
x_1 + x_2 \leq 7
\]
\[
x_1 \geq 0, \quad x_2 \geq 0
\]
\end{Ejem}


\begin{Ejem}[Pig Diet Problem]

Tabla nutricional:

\begin{center}
\begin{tabular}{|c|c|c|c|c|}
\hline
 & Corn & Silage & Alfalfa & Requerimiento diario \\
\hline
Carbs & 0.9 & 0.2 & 0.4 & 2 \\
Protein & 3 & 8 & 6 & 18 \\
Vitamins & 1 & 2 & 4 & 15 \\
Cost & 7 & 6 & 5 & \\
\hline
\end{tabular}
\end{center}
\end{Ejem}


\begin{Ejem}[Diet Model (G1 y G2)]

\begin{center}
\begin{tabular}{|c|c|c|c|c|}
\hline
 & Starch & Proteins & Vitamins & Cost (\$/kg) \\
\hline
G1 & 5 & 4 & 2 & 0.6 \\
G2 & 7 & 2 & 1 & 0.35 \\
\hline
\end{tabular}
\end{center}

Requerimientos diarios:
\[
\text{Starch} = 8, \quad \text{Proteins} = 15, \quad \text{Vitamins} = 3
\]

\end{Ejem}


\begin{Ejem}[Giapetto's Woodcarving]

Variables:
\[
x_1 = \text{soldiers}, \quad x_2 = \text{trains}
\]

Maximizar utilidad semanal considerando:

\begin{itemize}
\item Recursos de acabado y carpintería
\item 100 horas de acabado
\item 80 horas de carpintería
\item Máximo 40 soldados por semana
\end{itemize}

\end{Ejem}


\begin{Ejem}[General Linear Programming Problem]

Maximizar:
\[
\text{Max Profit} = 3x_1 + 4x_2
\]

Sujeto a:
\[
2x_1 + 3x_2 \leq 16
\]
\[
4x_1 + 2x_2 \leq 16
\]
\[
x_1, x_2 \geq 0
\]
\end{Ejem}


\begin{Ejem}[Furnco manufactures desks and chairs]

Let:
\[
x_1 = \text{number of desks produced}
\]
\[
x_2 = \text{number of chairs produced}
\]

Each desk contributes \$40 profit and each chair \$25.
Each desk uses 4 units of wood and each chair 3 units.
At most 20 units of wood are available.
At least twice as many chairs as desks must be produced.

The LP formulation is:

\[
\max \quad z = 40x_1 + 25x_2
\]

\[
\text{s.t.}
\]
\[
4x_1 + 3x_2 \le 20
\]
\[
x_2 \ge 2x_1
\]
\[
x_1,x_2 \ge 0
\]

Optimal solution:
\[
x_1=2,\quad x_2=4,\quad z=180.
\]
\end{Ejem}


\begin{Ejem}[Giapetto's Woodcarving]

Let:
\[
x_1 = \text{number of soldiers produced per week}
\]
\[
x_2 = \text{number of trains produced per week}
\]

\[
\max \quad z = 3x_1 + 2x_2
\]

\[
\text{s.t.}
\]
\[
2x_1 + x_2 \le 100
\]
\[
x_1 + x_2 \le 80
\]
\[
x_1 \le 40
\]
\[
x_1,x_2 \ge 0
\]

\end{Ejem}

\begin{Ejem}[Farmer Jones]

\textbf{Method 1: Acres as variables}

\[
x_1 = \text{acres of wheat}
\]
\[
x_2 = \text{acres of corn}
\]

\[
\max \quad z = 100x_1 + 30x_2
\]

\[
\text{s.t.}
\]
\[
x_1 + x_2 \le 7
\]
\[
10x_2 \ge 30
\]
\[
10x_1 + 4x_2 \le 40
\]
\[
x_1,x_2 \ge 0
\]

\textbf{Method 2: Bushels as variables}

\[
x_1 = \text{bushels of wheat}
\]
\[
x_2 = \text{bushels of corn}
\]

\[
\max \quad z = 4x_1 + 3x_2
\]

\[
\text{s.t.}
\]
\[
\frac{1}{25}x_1 + \frac{1}{10}x_2 \le 7
\]
\[
\frac{x_2}{10} \ge 3
\]
\[
\frac{1}{25}x_1 + \frac{1}{10}x_2 \le 40
\]
\[
x_1,x_2 \ge 0
\]

\end{Ejem}

\begin{Ejem}[Village Butcher Shop]

Let:
\[
x_1 = \text{pounds of ground beef}
\]
\[
x_2 = \text{pounds of ground pork}
\]

\[
\min \quad z = 80x_1 + 60x_2
\]

\[
\text{s.t.}
\]
\[
0.20x_1 + 0.32x_2 \le 0.25
\]
\[
x_1 + x_2 = 1
\]
\[
x_1,x_2 \ge 0
\]

\end{Ejem}


\begin{Ejem}[Psychologist Experiment]

Let:
\[
x = \text{number of mice}
\]
\[
y = \text{number of rats}
\]

\[
\max \quad z = x + y
\]

\[
\text{s.t.}
\]
\[
10x + 20y \le 800
\]
\[
20x + 10y \le 640
\]
\[
x,y \ge 0
\]

\end{Ejem}


